\documentclass{article}
\usepackage{amsmath}

\begin{document}

Test paragraph from SSIT section:

Each candidate wave has a certain likelihood $p$ of being interpreted as an IED by a rater, which we express on the logit scale $z=\log(p/(1-p))$; note that $z$ can be positive or negative. The rater's mental processes add noise (uncertainty) $n$ to $z$, producing a noisy impression, $z' = z+n$. The rater finally applies a threshold $\theta$ to reach a decision.

The model has two parameters for each rater: the noise/uncertainty level $\sigma$, which reflects a rater's skill in recognizing IEDs (an ideal rater would have a noise level of 0), and the threshold $\theta$, which represents the rater's personal preference as an over- or under-caller.

Estimated internal noise levels $\sigma$ and threshold $\theta$ for experts (blue), experienced non-experts (green), and novices (red) are shown in Figure 4A. As shown in Table 2, the novice group has larger noise/uncertainty $\sigma = 2.664$ [1.724, 8.577]) values compared to experienced non-experts $\sigma = 1.541$ [1.346, 1.838], and experts $\sigma = 1.262$ [1.061, 1.575].

\end{document}
